% ****************************************************************************************************
% Capítulo 4: Etiologia da Doença Mental
% ****************************************************************************************************
\chapter{Etiologia da Doença Mental}
\label{ch:etiologia}

% ****************************************************************************************************
\section{Onde Mora a Doença?}
\label{sec:onde-mora}
% ****************************************************************************************************

A pergunta pela etiologia --- a causa, a origem, o fundamento --- de uma doença mental parece, à primeira vista, pertencer ao domínio das ciências naturais. Identificar o agente causador, localizar a lesão, determinar o mecanismo fisiopatológico: este é o programa da medicina científica desde Virchow e sua patologia celular. Entretanto, como observa Jaspers (1913/1997), a psiquiatria ocupa posição peculiar entre as disciplinas médicas, pois seu objeto --- a experiência subjetiva alterada --- resiste às categorias explicativas que funcionam tão bem para tumores, infecções e fraturas.

Se perguntamos \enquote{onde está a pneumonia?}, a resposta é relativamente simples: nos pulmões, nas células alveolares inflamadas, no tecido infiltrado por neutrófilos. Mas se perguntamos \enquote{onde está a depressão?} --- onde, precisamente, ela \textit{mora} ---, a resposta já não é tão evidente. No cérebro? Em quais circuitos? Em quais neurotransmissores? Ou talvez na história de perdas acumuladas? Na infância marcada por negligência? No casamento que sufoca? No trabalho que aliena? Na sociedade que exclui? Na linguagem que nomeia o sofrimento de maneiras que o amplificam ou atenuam?

Este capítulo propõe uma investigação sobre o \foreign{locus} etiológico das condições que chamamos \enquote{doenças mentais}. Argumenta-se que a resposta a esta pergunta --- aparentemente técnica, aparentemente neutra --- carrega implicações profundas sobre como concebemos o sofrimento psíquico, como o tratamos, e como distribuímos responsabilidades entre o indivíduo que sofre e o mundo que o constitui.

A tese que se desenvolverá é simples, talvez óbvia demais para parecer provocativa: \textit{a etiologia de uma doença mental não pode ser localizada exclusivamente no interior do indivíduo que a manifesta}.

% ****************************************************************************************************
\section{A Temporalidade do Adoecer Mental}
\label{sec:temporalidade}
% ****************************************************************************************************

\subsection{A Doença Como Evento ou Como Processo?}

A medicina somática frequentemente opera com modelo de doença como \textit{evento} --- algo que \textit{acontece} ao organismo em momento determinado. O infarto ocorre às 14h37 de uma terça-feira; a fratura acontece no instante do impacto; a infecção inicia quando o patógeno penetra as barreiras do hospedeiro.

Mas quando começa uma depressão? Quando começa uma esquizofrenia?

Fuchs (2012), em seu trabalho sobre temporalidade e psicopatologia, observa que os transtornos mentais caracterizam-se por estrutura temporal radicalmente distinta. Não há o \textit{momento} em que alguém \enquote{contrai} depressão como se contrai gripe. Antes, o que observamos é processo gradual de transformação --- o que Minkowski (1927/1970) denominou alteração no \enquote{tempo vivido} (\foreign{temps vécu}) --- que frequentemente só se torna visível retrospectivamente.

Esta característica temporal tem consequências etiológicas fundamentais. Se a doença não \textit{começa} em momento identificável, mas \textit{emerge} de processo prolongado, então a busca por \textit{a} causa --- singular, pontual, determinante --- revela-se inadequada. Como observa Kendler (2012), a causalidade em psiquiatria é \enquote{matizada} (\foreign{dappled}): múltiplos fatores interagem ao longo do tempo, nenhum suficiente isoladamente, nenhum necessário absolutamente.

\subsection{A Doença Tem História}

Ricoeur (1990/1992), em sua análise da identidade narrativa, propõe que o \textit{si mesmo} não é substância fixa, mas configuração temporal --- somos as histórias que contamos sobre nós mesmos. Se o sujeito é narrativa, então sua doença também é narrativa --- não estado atemporal, mas \textit{história} com enredo, personagens, cenários.

Considere-se um caso paradigmático: uma mulher de 45 anos apresenta episódio depressivo grave. O psiquiatra pode documentar os sintomas, pode localizar alterações em circuitos neurais, pode prescrever antidepressivo. Mas \textit{esta} depressão, \textit{desta} mulher, não é apenas conjunto de sintomas. É também: o pai alcoólatra que a humilhava na infância; o primeiro casamento marcado por violência; a demissão após vinte anos de dedicação; os filhos que se afastaram; o corpo que envelhece em sociedade que valoriza juventude; a solidão em cidade que não reconhece.

Onde está a etiologia \textit{desta} depressão? No déficit de serotonina? Na história de trauma? Na estrutura econômica que permite demissões arbitrárias? Nos valores culturais que depreciam mulheres maduras? A pergunta, formulada assim, já contém sua resposta: em todos estes lugares, e em nenhum deles exclusivamente.

\subsection{Temporalidade e Causalidade Circular}

A relação entre fatores etiológicos e manifestação sintomática raramente é linear. Como observa Bateson (1972), o que encontramos em fenômenos complexos é \textit{causalidade circular} --- A causa B que causa C que retroage sobre A, em loops que se auto-mantêm e auto-amplificam.

O isolamento social pode causar depressão; a depressão causa mais isolamento; o isolamento aprofunda a depressão. O trauma gera desconfiança; a desconfiança impede vínculos reparadores; a ausência de vínculos perpetua os efeitos do trauma.

Nestes circuitos, perguntar pela causa \textit{primeira} é exercício frequentemente infrutífero.

% ****************************************************************************************************
\section{A Contextualidade do Adoecer Mental}
\label{sec:contextualidade}
% ****************************************************************************************************

\subsection{O Contexto Como Fundo Necessário}

Merleau-Ponty (1945/2006), em sua fenomenologia da percepção, demonstra que toda figura só é figura contra um fundo --- não existe percepção de objeto isolado, mas sempre de objeto \textit{em} campo perceptivo que o circunda e o constitui.

Se aplicamos esta insight à psicopatologia, reconhecemos que o sintoma --- a figura que captura atenção clínica --- só é sintoma contra fundo contextual que lhe confere significado. A tristeza após perda de ente querido não é o mesmo fenômeno que a tristeza sem motivo aparente, ainda que a descrição fenomenológica possa ser idêntica. O contexto não é mero \enquote{cenário} onde a doença se desenrola; o contexto \textit{constitui} a doença.

Kleinman (1988), em seus estudos sobre narrativas de doença, distingue \foreign{disease} (a disfunção biológica objetivável) de \foreign{illness} (a experiência vivida do adoecimento, culturalmente mediada). A \enquote{mesma} depressão é vivida diferentemente pelo executivo paulistano e pelo agricultor nordestino, pela adolescente e pela idosa, pelo heterossexual e pelo homossexual em sociedade homofóbica.

\subsection{Contextos Patogênicos}

Se o contexto constitui a doença, então certos contextos são \textit{patogênicos} --- não no sentido de conterem \enquote{germes} psíquicos, mas no sentido de criarem condições nas quais o adoecimento torna-se mais provável, mais grave, mais persistente.

A literatura epidemiológica documenta abundantemente estes contextos: pobreza, desigualdade, discriminação, violência, migração forçada, precarização do trabalho. Como observa Wilkinson e Pickett (2009), sociedades mais desiguais apresentam prevalência maior de transtornos mentais --- não apenas entre os pobres, mas na população geral. A desigualdade é tóxico social que afeta o tecido relacional como um todo.

Mas se contextos podem ser patogênicos, então a etiologia da doença mental que emerge nestes contextos não pode ser localizada exclusivamente no indivíduo que adoece.

\subsection{A Pergunta Incômoda}

Chegamos aqui a pergunta que a psiquiatria tradicional prefere não formular de maneira tão direta: \textit{se o contexto é patogênico, por que tratamos apenas o indivíduo?}

A resposta habitual --- que não podemos mudar a sociedade, mas podemos tratar o paciente --- é pragmaticamente compreensível, mas epistemologicamente insustentável. É como se, diante de epidemia de cólera, a medicina se limitasse a reidratar os doentes sem questionar a qualidade da água.

Como observa Rose (2019), a psiquiatria contemporânea opera com o que ele denomina \enquote{estilo de pensamento neurobiológico} --- tendência a localizar a etiologia dos transtornos mentais em disfunções cerebrais individuais, obscurecendo determinantes sociais. Este estilo de pensamento não é politicamente neutro; ao individualizar o sofrimento, despolitiza-o.

% ****************************************************************************************************
\section{O Locus Etiológico: Onde Está a Causa?}
\label{sec:locus-etiologico}
% ****************************************************************************************************

\subsection{A Etiologia no Eu: O Modelo Biomédico}

O modelo biomédico tradicional localiza a etiologia da doença mental no organismo individual --- mais especificamente, no cérebro. Transtornos mentais seriam \enquote{doenças do cérebro} causadas por disfunções em circuitos neurais específicos.

Este modelo possui méritos inegáveis. A demonstração de correlatos neurobiológicos em condições como esquizofrenia e transtorno bipolar legitimou o sofrimento mental como fenômeno \enquote{real}. Entretanto, como observa Kirmayer e Gold (2012), o modelo opera com redução epistêmica problemática --- a suposição de que níveis superiores de organização (experiência, significado, relação) podem ser integralmente explicados por níveis inferiores (neurônios, sinapses, moléculas).

A pergunta crítica é: \textit{mesmo admitindo que todo fenômeno mental tenha correlatos neurobiológicos, segue-se que a etiologia está no cérebro?} A correlação não é causação.

\subsection{A Etiologia no Outro: O Modelo Relacional}

Uma tradição alternativa --- que inclui psicanálise das relações objetais, teoria do apego, psicologia interpessoal --- localiza a etiologia não no indivíduo isolado, mas nas \textit{relações} que o constituíram.

Winnicott (1960/1983), com sua célebre observação de que \enquote{não existe bebê} --- apenas díade mãe-bebê ---, fundamenta esta perspectiva. O self não é dado originário; o self \textit{emerge} das relações. Se as relações primárias são suficientemente boas, o self desenvolve-se integrado. Se são traumáticas, o self desenvolve-se fragmentado.

Bowlby (1969/1990), fundador da teoria do apego, documenta empiricamente esta dependência constitutiva. Padrões de apego formados na infância predizem funcionamento psicológico ao longo da vida. Onde está a etiologia da ansiedade do adulto ansiosamente apegado? Nele, certamente --- é ele quem a experimenta. Mas também no \textit{outro} que falhou em oferecer base segura.

\subsection{A Etiologia na Sociedade: O Modelo Crítico}

Uma terceira perspectiva --- associada a Foucault, à antipsiquiatria, aos estudos críticos em saúde mental --- localiza a etiologia não no indivíduo nem nas relações diádicas, mas nas \textit{estruturas sociais} que produzem e distribuem sofrimento de maneira desigual.

Foucault (1961/2006), em sua arqueologia da loucura, demonstra que as categorias psiquiátricas não são descobertas neutras, mas construções históricas inseparáveis de dispositivos de poder. A loucura não é simplesmente \textit{encontrada} pela psiquiatria; é, em certo sentido, \textit{produzida} por ela.

Considere-se a epidemia contemporânea de \enquote{transtornos de ansiedade}. Seria coincidência que esta epidemia ocorra precisamente em sociedades caracterizadas por precarização do trabalho, erosão de vínculos comunitários, competição generalizada, exposição constante a notícias ameaçadoras? Ou estamos medicalizando --- e consequentemente despolitizando --- respostas razoáveis a condições sociais objetivamente ansiogênicas?

\subsection{A Etiologia na Linguagem: O Modelo Discursivo}

Uma quarta perspectiva --- derivada da psicanálise lacaniana, da filosofia da linguagem, dos estudos de gênero --- localiza a etiologia na \textit{linguagem} e nos \textit{discursos} que constituem a subjetividade.

Lacan (1966/2006) propõe que \enquote{o inconsciente é estruturado como linguagem}. Não somos sujeitos que \textit{usam} linguagem; somos sujeitos \textit{constituídos} pela linguagem, atravessados por significantes que nos precedem e nos determinam.

Butler (1990/2003), em sua teoria da performatividade de gênero, demonstra que categorias como \enquote{masculino} e \enquote{feminino} não descrevem essências naturais, mas são performativamente constituídas. O sofrimento que resulta desta constituição discursiva --- a disforia de gênero, a vergonha do corpo --- é real. Mas onde está sua etiologia? No indivíduo que sofre? Ou nos discursos normativos que definem certos corpos como patológicos?

% ****************************************************************************************************
\section{Implicações da Deslocalização Etiológica}
\label{sec:implicacoes}
% ****************************************************************************************************

\subsection{Implicações para o Diagnóstico}

Se a etiologia não está (apenas) no indivíduo, então o diagnóstico não pode (apenas) catalogar sintomas individuais. A avaliação clínica deve incluir o que Kleinman (1988) denomina \enquote{modelo explicativo} do paciente, mas também avaliação do contexto: Quais relações sustentam ou adoecem este paciente? Quais condições sociais contribuem para seu sofrimento? Quais discursos o constituíram como \enquote{doente} ou \enquote{desviante}?

\subsection{Implicações para o Tratamento}

Se a etiologia está distribuída entre indivíduo, relações, sociedade e discursos, então o tratamento também deve ser distribuído. A intervenção farmacológica pode ser necessária, mas não é suficiente. A psicoterapia individual pode ser útil, mas não basta se as relações do paciente permanecem tóxicas. A mudança individual pode ser desejável, mas não é justa se a sociedade que adoece permanece intocada.

\subsection{Implicações para a Responsabilidade}

A pergunta pela etiologia é também pergunta pela \textit{responsabilidade}. Se a doença está no indivíduo, então é ele quem deve mudar. Se a doença está nas relações, na sociedade, nos discursos, então a responsabilidade é compartilhada --- e exigir que apenas o indivíduo mude é, além de ineficaz, injusto.

% ****************************************************************************************************
\section{A Etiologia Como Questão Aberta}
\label{sec:etiologia-aberta}
% ****************************************************************************************************

Iniciamos este capítulo com pergunta aparentemente simples: onde está a causa da doença mental? Percorremos possíveis respostas --- no Eu, no Outro, na Sociedade, na Linguagem --- e encontramos, em cada caso, verdade parcial e limitação.

A resposta honesta, talvez, seja que a etiologia da doença mental é \textit{distribuída} --- não localizável em ponto único, mas dispersa através de múltiplos níveis de organização que interagem de maneiras complexas. O cérebro importa, certamente --- mas o cérebro existe em corpo, o corpo existe em relações, as relações existem em sociedade, a sociedade existe em história, a história existe em linguagem.

Esta conclusão pode parecer frustrante para quem busca respostas definitivas. Mas é, talvez, também libertadora. Se a etiologia está em múltiplos lugares, então a intervenção pode ocorrer em múltiplos pontos. Não precisamos esperar que a neurociência \enquote{descubra} a causa da depressão para agir sobre condições sociais que a geram.

Pois se a doença mental tem sua origem, ao menos parcialmente, fora do indivíduo que a manifesta, então a tarefa do clínico não é apenas tratar o paciente, mas também testemunhar a injustiça que o adoece. E testemunhar é já uma forma de cura.

% ****************************************************************************************************
% Referências do Capítulo
% ****************************************************************************************************
\vspace{2em}
\begin{small}
\noindent\textsc{Referências}

\vspace{1em}
\noindent Bateson, G. (1972). \textit{Steps to an ecology of mind}. University of Chicago Press.

\noindent Bowlby, J. (1990). \textit{Apego e perda: Vol. 1. Apego}. (Orig. 1969). Martins Fontes.

\noindent Butler, J. (2003). \textit{Problemas de gênero}. (Orig. 1990). Civilização Brasileira.

\noindent Foucault, M. (2006). \textit{História da loucura}. (Orig. 1961). Perspectiva.

\noindent Fuchs, T. (2012). Temporality and psychopathology. \textit{Phenomenology and the Cognitive Sciences}, 12(1), 75--104.

\noindent Kendler, K. S. (2012). The dappled nature of causes of psychiatric illness. \textit{Molecular Psychiatry}, 17(4), 377--388.

\noindent Kirmayer, L. J., \& Gold, I. (2012). Re-socializing psychiatry. In \textit{Critical neuroscience}. Wiley-Blackwell.

\noindent Kleinman, A. (1988). \textit{The illness narratives}. Basic Books.

\noindent Rose, N. (2019). \textit{Our psychiatric future}. Polity.

\noindent Wilkinson, R. G., \& Pickett, K. (2009). \textit{The spirit level}. Allen Lane.

\noindent Winnicott, D. W. (1983). O ambiente e os processos de maturação. (Orig. 1960). Artes Médicas.
\end{small}
